\cvsection{Expériences professionnelles}

\begin{cventries}
  \cventry
    {Stagiaire ingénieur de recherche en IA}
    {INRIA}
    {Bordeaux, France}
    {Mars 2026 -- Septembre 2026}
    {
      \begin{cvitems}
        \item {Définition d’un protocole expérimental pour comparer Qwen, Gemma et GLM selon la qualité des réponses, la latence, le TTFT et la taille du contexte.}
        \item {Évaluation de plusieurs quantifications pour étudier les compromis entre qualité et performances selon les environnements d’exécution.}
        \item {Installation, configuration et serving de LLM avec vLLM sur une infrastructure HPC équipée de GPU NVIDIA, avec utilisation de Slurm.}
        \item {Analyse des résultats expérimentaux et rédaction en cours d’un article scientifique comparant les modèles pour notre cas d’usage.}
        \item {Conception et développement en autonomie d’une application et d’un agent IA en Python, avec services FastAPI, interface React/Next.js et pipeline RAG.}
        \item {Conteneurisation de composants avec Docker et gestion d’une application utilisant des modèles locaux via LM Studio sur Mac Studio.}
      \end{cvitems}
    }

  \cventry
    {Stagiaire ingénieur de recherche en apprentissage par renforcement}
    {Université Karunya}
    {Coimbatore, Inde}
    {Juin 2025 -- Sept. 2025}
    {
      \begin{cvitems}
        \item {Développement d’un modèle d’apprentissage par renforcement et de son environnement d’entraînement pour la navigation autonome d’un robot hospitalier.}
        \item {Déploiement et optimisation du logiciel et du modèle sur Raspberry Pi pour des tests en conditions réelles.}
      \end{cvitems}
    }
\end{cventries}
